% Part: proof-theory
% Chapter: cut-elimination
% Section: introduction

\documentclass[../../../include/open-logic-section]{subfiles}

\begin{document}

\olfileid{pt}{cut}{int}

\olsection{سريزه}

د \Cut{} قاعده دا ده:
\[
\Axiom$\Gamma \fCenter \Delta, !A$
\Axiom$!A, \Pi \fCenter \Lambda$
\RightLabel{\Cut}
\BinaryInf$\Gamma, \Pi \fCenter \Delta, \Lambda$
\DisplayProof
\]
(!!{formula}~$!A$ ته \emph{د پرېکون فارمول} وايي۔)

\Cut{} د جېنڅن په اصلي سېکوېنټ حساب~\Log{LK} کښې شامله وه۔
د جېنڅن \emph{هاوپټ‌ساتڅ} («اصلي قضيه») وايي چې په~\Log{LK} کښې
هر !!{provable} سېکوېنټ د \Cut{} له قاعدې کار اخستلو پرته هم
!!{provable} دے؛ يعنې \Cut{} د \Log{LK} له ثبوتونو څخه
«لرې کېدے شي»۔ زمونږ نظامونه \Log{G1c} او \Log{G3c}، \Cut نۀ لري،
خو د هغو دپاره هم پوښتنه کولے شو: آيا \Cut{} له هغو !!{proof}s
څخه لرې کېدے شي چې \Cut{} او د \Log{G1c} نور قواعد، يا د
\Log{G3c} قواعد، کاروي؟ په زمونږ ټاکل شوې اصطلاح کښې هم‌ارزښته
پوښتنه دا ده: آيا \Cut{} د \Log{G1c}، يا \Log{G3c}، يوه د منلو
وړ قاعده ده؟

د پرېکون لرې کولو قضيه ښايي د ثبوت پوهنې تر ټولو مهمه او
مشهوره پايله وي۔ تر هر څه مخکې ښيي چې \Log{G1c} او \Log{G3c}
هغې قاعدې ته اړ نۀ دي چې په لومړي نظر درته لازمي ښکاري۔ کله
چې واقعي ثبوتونه رسمي کوئ، د مرستندويه قضيو له استعمال سره د
کار طريقه غواړئ۔ رياضيپوهان پايلې ثابتوي او بيا د نورو پايلو
په ثبوتونو کښې ورته رجوع کوي۔ نو لومړے د پايلې $L$، يعنې
مرستندويې قضيې، ثبوت د سېکوېنټ $\Sequent L$ د !!{proof} په توګه
رسمي کولے شئ۔ بيا د دويمې پايلې~$R$ ثبوت، چې~$L$ کاروي، د
سېکوېنټ $L \Sequent
R$ د !!{proof} په توګه رسمي کوئ۔ څنګه پوهېږئ چې~$R$ ثبوت لري؛
يعنې يوازې د $\Sequent R$ يو !!a{proof} شته؟ د دواړو !!{proof}s
د يوځاے کولو طريقه غواړئ، او د \Cut{} قاعده دا اجازه درکوي۔

يادونه مو کړې وه چې \Log{G1c} او \Log{G3c} بشپړ دي، نو ټول
هغه سېکوېنټونه ثابتوي چې د ثابتېدو تمه ئې کوئ؛ يعنې ټول معتبر
سېکوېنټونه۔ دا نېغ په نېغه ثابتېدے شي۔ خو د جېنڅن د کار په
وخت کښې يوازې بديهياتي اشتقاقي نظامونه بشپړ معلوم وو۔ د
\Log{LK} د بشپړتيا د ښودلو دپاره جېنڅن د بديهياتي نظام
!!{derivation}s د \Log{LK} په !!{proof}s بدل کړل۔ په دې بدلولو
کښې \Cut{} اساسي نقش درلود۔ د پرېکون لرې کول يوازې د کلاسيکي
منطق مهمه پايله نۀ ده: ډېر نور منطقونه هم سېکوېنټ حسابونه لري۔
د ځينو منطقونو د سېکوېنټ حسابونو د بشپړتيا نېغ ثبوت سخت يا
ناممکن وي، نو د نورو نظامونو څخه، د \Cut په کارولو، بدلول ئې
د بشپړتيا يوازېنۍ طريقه وي۔ بيا په ثبوت پوهنه د پرېکون لرې
کول ثابتول اړين دي، څو وښيي چې \Cut{} اضافي ده او له \Cut{}
پرته نظام بشپړ دے۔

د پرېکون لرې کول ځکه هم مهم دي چې د نورو، په خپل ځاے په زړه
پورې، پايلو د تر لاسه کولو دپاره کارېدے شي۔ دوه بېلګې چې
ګورو ئې د کرېګ د منځګړيتوب مرستندويه قضيه او د هېربراند قضيه دي۔

د پرېکون لرې کولو د ثبوت طريقې عموماً ښيي چې له \Cut{} څخه
کار اخستونکے !!a{proof} څنګه په بې‌پرېکونه ثبوت بدلېدے شي۔ دا
بدلونونه اکثر «ځايي» وي؛ يعنې د يو !!
a{proof} يوه برخه، عموماً په \Cut{} استنتاج ختمېدونکے فرعي
!!{proof}، د هماغه سېکوېنټ په بل فرعي !!{proof} بدلوو چې پکښې
د \Cut{} استنتاج لرې يا په نورو بدل شوے وي۔ داسې استدلالونه
له \Cut{} سره !!{proof}s په بې‌پرېکونه ثبوتونو د بدلولو ګام
په ګام الګوريتمونه ورکوي۔ دا د پرېکون لرې کولو د طريقو په
استعمال کښې نور معلومات راکوي۔ مثلاً د مدل پوهنې په طريقه
د منځګړيتوب د مرستندويې قضيې داسې ثبوتونه شته: که $!A \lif !B$
معتبر وي، نو منځګړی شته (اوس دا پرېږدئ چې منځګړی څه ته وايي)۔
د ثبوت پوهنې استدلال د پرېکون لرې کولو په وسيله داسې الګوريتم
ورکوي چې د $!A
\lif !B$ يو !!a{proof} اخلي او منځګړی راګرځوي۔ د مدل پوهنې
له استدلاله بېل، د ثبوت پوهنې استدلال \emph{رغنده} دے۔

د پرېکون لرې کولو د قضيې د ثابتولو دوه مشهورې طريقې دي۔
يوه، چې جېنڅن ته ورګرځي، ښيي چې د ثبوت تر ټولو پاسني پرېکونونه
لرې کولے شئ، او بيا ئې تر هغه لرې کوي چې يو هم پاتې نۀ شي۔
بله، چې لومړے شوتې او ټېټ وړاندې کړې ده، په يوه !!a{proof}
کښې تر ټولو پېچلو پرېکونونو ته ګوري او هغه پرېکونونه په وار
لرې کوي چې د دې معنا له مخې اعظمي دي: هېڅ بل پرېکون تر هغو
ژور د پرېکون !!{formula} نۀ لري۔ هره طريقه خپلې ګټې او ستونزې
لري۔ د ځينو نظامونو دپاره يوه طريقه کار کوي او بله سختېږي يا
ناممکنه وي۔ نو دواړه طريقې بيانوو۔ د~\Log{G3c} دپاره د
پرېکون لرې کول ثابتوو۔ په حقيقت کښې ښيو چې يوازې \Cut{} نۀ،
بلکې \emph{د ګډ سياق پرېکون}~\CutCS لرې کېدے شي:
\[
\Axiom$\Gamma \fCenter \Delta, !A$
\Axiom$!A, \Gamma \fCenter \Delta$
\RightLabel{\CutCS}
\BinaryInf$\Gamma \fCenter \Delta$
\DisplayProof
\]
\textbf{سرچينه‌يي سمون OLCMP-123:} د دې رسم نښه د ګډ سياق
پرېکون ده؛ دواړه مقدمې عين سياق لري او پايله ئې يو نقل ساتي۔
که سېکوېنټ حساب کمزوري کول او انقباض مني، که د \Log{G1c} په
شان اصلي قواعد وي يا د \Log{G3c} په شان د منلو وړ قواعد، نو
\Cut{}، \CutCS تمثيلولے شي او معکوس ئې هم شته۔ د \Cut پر
ځاے \CutCS{} غوره کوو، ځکه د پرېکون لرې کولو لومړۍ طريقه
د~\CutCS دپاره اسانه بيانېږي او دويمه طريقه يوازې د~\CutCS
دپاره کار کوي۔

\end{document}
